\documentclass[11pt,a4paper]{article}

\usepackage[utf8]{inputenc}
\usepackage[T1]{fontenc}
\usepackage{lmodern}
\usepackage[margin=2.4cm]{geometry}
\usepackage{array}
\usepackage{longtable}
\usepackage{hyperref}
\usepackage{xcolor}

\hypersetup{
  colorlinks=true,
  linkcolor=blue,
  urlcolor=blue,
  pdftitle={Claude Code mit OpenCode Go als Backend verwenden},
  pdfauthor={claude-o-go}
}

\setlength{\parindent}{0pt}
\setlength{\parskip}{0.7em}
\renewcommand{\arraystretch}{1.2}

\newcommand{\code}[1]{\texttt{#1}}

\title{\Huge Claude Code mit OpenCode Go\\[0.4em]\Large Ein praktisches Backend-Tutorial fuer \code{claude-o-go}}
\author{\code{claude-o-go}}
\date{\today}

\begin{document}
\maketitle
\tableofcontents
\newpage

\section{Ziel}

Dieses Tutorial erklaert den Aufbau dieses Repositories und warum er
funktioniert. Ziel ist, die Claude Code CLI weiter zu verwenden, aber
Modellanfragen an OpenCode Go statt an Anthropic Claude zu senden.

Claude Code bleibt der Harness: Terminal UI, Tool Calling, Dateioperationen,
Permissions, Prompt-Verarbeitung und Agent Loop. OpenCode Go ist das Backend,
das die eigentliche Modellanfrage erhaelt.

\section{Was vorhanden ist}

Dieses Verzeichnis enthaelt:

\begin{itemize}
  \item \code{.env}: dein OpenCode Go API Key.
  \item \code{claude-o-go}: das Launcher-Skript, das du statt \code{claude} startest.
  \item \code{opencode-go-claude-proxy.mjs}: der lokale Proxy zwischen Claude Code und OpenCode Go.
  \item \code{opencode-go-claude-example.mjs}: ein direkter API-Test ohne Claude Code.
\end{itemize}

Der normale Start:

\begin{verbatim}
cd /home/user/Desktop/opencode-go-w-claude
./claude-o-go
\end{verbatim}

Ein einzelner Prompt:

\begin{verbatim}
./claude-o-go -p "Reply with exactly: OK"
\end{verbatim}

\section{\code{claude-o-go} von ueberall starten}

Du musst nicht jedes Mal in dieses Repository wechseln und es nicht in jedes
Projekt kopieren. Installiere es einmal und verwende \code{claude-o-go} aus
jedem Projektverzeichnis.

\subsection{Einmalige Einrichtung}

Launcher ausfuehrbar machen, falls noetig:

\begin{verbatim}
chmod +x /home/user/Desktop/opencode-go-w-claude/claude-o-go
\end{verbatim}

Symlink in einem Verzeichnis auf deinem \code{PATH} anlegen:

\begin{verbatim}
mkdir -p ~/.local/bin
ln -s /home/user/Desktop/opencode-go-w-claude/claude-o-go \
  ~/.local/bin/claude-o-go
\end{verbatim}

Sicherstellen, dass \code{\textasciitilde{}/.local/bin} auf dem \code{PATH} liegt:

\begin{verbatim}
export PATH="$HOME/.local/bin:$PATH"
\end{verbatim}

Shell neu laden:

\begin{verbatim}
source ~/.bashrc   # oder: source ~/.zshrc
\end{verbatim}

\subsection{Warum das funktioniert}

Der Launcher loest seinen eigenen Symlink mit \code{readlink -f} auf. Dadurch
findet er das echte Installationsverzeichnis, in dem \code{.env} und
\code{opencode-go-claude-proxy.mjs} liegen. Claude Code selbst startet wieder
in dem Verzeichnis, aus dem du \code{claude-o-go} aufgerufen hast, sodass dein
aktuelles Projekt das Arbeitsverzeichnis bleibt.

\subsection{Optionale Shell-Konfiguration}

Wenn du den Key nicht in \code{.env} halten willst, exportiere ihn einmal in
deinem Shell-Profil:

\begin{verbatim}
export OPENCODE_GO_API_KEY="sk-..."
export OPENCODE_GO_MODEL="qwen3.7-plus"
\end{verbatim}

Der Launcher nimmt \code{OPENCODE\_GO\_API\_KEY} aus der Umgebung und
ueberspringt dann \code{.env}. Das Modell ist nicht mehr im Launcher
hartcodiert; setze \code{OPENCODE\_GO\_MODEL} im Shell-Profil, damit es ueberall
gilt.

\subsection{Taegliche Nutzung}

Aus einem beliebigen Projekt:

\begin{verbatim}
cd /path/to/any/project
claude-o-go
\end{verbatim}

Ein einzelner Prompt:

\begin{verbatim}
claude-o-go -p "Explain this repository in one paragraph."
\end{verbatim}

Modell fuer einen Lauf wechseln:

\begin{verbatim}
OPENCODE_GO_MODEL=minimax-m3 claude-o-go
\end{verbatim}

\section{Warum ein Proxy noetig ist}

OpenCode Go stellt mehrere Modelle ueber eine Anthropic-kompatible Messages API
bereit. Es ist eine Routing- und Abrechnungsschicht, kein Modell-Host. Jeder
Request wird an die API des jeweiligen Modellanbieters weitergeleitet.

Beispiele:

\begin{verbatim}
qwen3.7-plus
minimax-m3
qwen3.7-max
\end{verbatim}

Claude Code validiert Modellnamen lokal, bevor ein Request abgeschickt wird.
Ein direkter Start mit einem OpenCode-Go-Modell wird deshalb abgelehnt:

\begin{verbatim}
ANTHROPIC_BASE_URL=https://opencode.ai/zen/go/v1 \
ANTHROPIC_API_KEY=... \
claude --model qwen3.7-plus
\end{verbatim}

Der Proxy loest das, indem Claude Code lokal ein normales Claude-Alias wie
\code{sonnet} sieht. Der Proxy schreibt den ausgehenden Request dann auf das
OpenCode-Go-Modell um.

\section{Exkurs: Localhost, Ports und Prozesse}

Networking beginnt mit einer einfachen Frage: wohin soll ein Client verbinden?

\begin{itemize}
  \item \code{127.0.0.1} bedeutet localhost. Die Adresse zeigt auf dieselbe Maschine.
  \item Ein Port waehlt einen konkreten lauschenden Prozess aus. Der Standard-Port ist \code{4141}.
  \item Ein Server bindet sich an Adresse und Port. Der Proxy bindet sich an \code{127.0.0.1:4141}.
  \item Ein Client verbindet sich zu Adresse und Port. Claude Code verbindet sich zur lokalen Proxy-URL:
\end{itemize}

\begin{verbatim}
http://127.0.0.1:4141
\end{verbatim}

\begin{verbatim}
Claude Code Prozess
  verbindet zu 127.0.0.1:4141

Node Proxy Prozess
  lauscht auf 127.0.0.1:4141
  verbindet nach aussen zu opencode.ai:443
\end{verbatim}

Der Proxy lauscht auf localhost, nicht auf \code{0.0.0.0}. Das ist wichtig:
\code{127.0.0.1} ist nur von deiner eigenen Maschine erreichbar.
\code{0.0.0.0} wuerde den Listener auf allen Netzwerkinterfaces anbieten.

Nuetzliche Admin-Checks:

\begin{verbatim}
ss -ltnp | grep 4141
lsof -iTCP:4141 -sTCP:LISTEN
\end{verbatim}

\section{Exkurs: HTTP-Proxy vs Backend}

Claude Code ist der Client. OpenCode Go ist das Upstream-Backend. Das lokale
Node-Skript ist ein Proxy.

Ein Proxy nimmt einen Request von einer Seite an und sendet einen verwandten
Request an eine andere Seite. Er kann den groessten Teil unveraendert lassen,
einzelne Felder umschreiben und danach die Antwort zurueckreichen.

In diesem Projekt macht der Proxy drei wichtige Dinge:

\begin{enumerate}
  \item Er nimmt Claude Codes Anthropic-Request unter \code{/v1/messages} an.
  \item Er schreibt \code{body.model} von Claudes lokalem Alias auf das OpenCode-Go-Modell um.
  \item Er leitet den Request an OpenCode Go weiter und streamt die Antwort zurueck.
\end{enumerate}

Der Proxy ist kein Modell-Host. Er erzeugt keine Tokens. Er ist Glue-Code
zwischen zwei APIs, die fast, aber nicht perfekt, kompatibel sind.

\section{Modell-Kompatibilitaet}

Jeder Upstream-Anbieter implementiert seinen Anthropic-kompatiblen Endpoint
selbst. Deshalb variiert die Kompatibilitaet. Claude Code sendet einen
vollstaendigen Anthropic-Payload mit Tools, Cache Control, Tool Choice und
Anthropic-Beta-Headern.

\begin{longtable}{|p{0.24\linewidth}|p{0.22\linewidth}|p{0.43\linewidth}|}
\hline
\textbf{Modell-ID} & \textbf{Anbieter} & \textbf{Funktioniert mit Claude Code Harness?} \\
\hline
\code{qwen3.7-plus} & Alibaba (Qwen) & Ja \\
\hline
\code{minimax-m3} & MiniMax & Ja \\
\hline
\code{glm-5.2} & Z.ai (Zhipu AI) & Nein; invalid API parameter \\
\hline
\code{glm-5.1} & Z.ai (Zhipu AI) & Nein; wie \code{glm-5.2} \\
\hline
\code{kimi-k2.7-code} & Moonshot AI & Nein; lehnt Tool Function Names ab \\
\hline
\code{deepseek-v4-pro} & DeepSeek & Nein; erwartet OpenAI Tool Shape \\
\hline
\end{longtable}

Modelle desselben Anbieters verhalten sich meist wie das getestete Modell.
Ungepruefte Modelle sollten vor produktiver Nutzung getestet werden.

Praktische Empfehlung:

\begin{verbatim}
export OPENCODE_GO_MODEL="qwen3.7-plus"   # oder minimax-m3
\end{verbatim}

\section{Request Flow}

Wenn du Folgendes startest:

\begin{verbatim}
./claude-o-go -p "Reply with exactly: OK"
\end{verbatim}

passiert Folgendes:

\begin{enumerate}
  \item \code{claude-o-go} liest den OpenCode Go API Key aus \code{.env}.
  \item Es startet den lokalen Proxy auf \code{http://127.0.0.1:4141}.
  \item Es setzt \code{ANTHROPIC\_BASE\_URL=http://127.0.0.1:4141}.
  \item Es startet Claude Code mit \code{claude --bare --model sonnet}.
  \item Claude Code sendet einen Request an \code{/v1/messages}.
  \item Der Proxy aendert das Modell im Request Body auf das OpenCode-Go-Modell.
  \item Der Proxy leitet an \code{https://opencode.ai/zen/go/v1/messages} weiter.
  \item OpenCode Go liefert die Modellantwort.
  \item Der Proxy streamt die Antwort zurueck an Claude Code.
\end{enumerate}

\section{Exkurs: Streaming, SSE und Backpressure}

LLM-Antworten werden normalerweise gestreamt. Das Backend wartet nicht, bis die
komplette Antwort fertig ist, sondern sendet kleine Chunks, sobald sie
verfuegbar sind.

Ein haeufiges Web-Muster dafuer ist SSE: Server-Sent Events. Ein SSE-Stream ist
weiterhin HTTP, aber der Response Body bleibt offen und sendet ueber Zeit
Events:

\begin{verbatim}
data: {"type":"content_block_delta","text":"hello"}

data: {"type":"content_block_delta","text":" world"}

data: [DONE]
\end{verbatim}

Backpressure bedeutet, dass der Empfaenger Daten nicht immer so schnell
verarbeiten kann, wie der Sender sie schreibt. In Node gibt
\code{response.write(chunk)} \code{false} zurueck, wenn der Ausgabepuffer voll
ist. Korrekte Proxy-Logik wartet dann auf \code{drain}, bevor sie
weiterschreibt. Das verhindert Speicherwachstum und instabile lange Streams.

In diesem Projekt:

\begin{itemize}
  \item Claude Code liest vom lokalen Proxy.
  \item Der Proxy liest von OpenCode Go.
  \item Der Proxy schreibt Chunks an Claude Code und beachtet Node-Stream-Backpressure.
  \item Wenn Claude Code die lokale Verbindung schliesst, bricht der Proxy den Upstream-Request ab.
\end{itemize}

Stream-Handling ist empfindlicher als ein einfacher JSON-Request. Ein normaler
Request hat eine Antwort. Ein Stream ist ein laufendes Gespraech zwischen
Sockets.

\section{Exkurs: Haeufige Netzwerkfehler}

Diese Fehlernamen gehoeren zum praktischen Admin-Wortschatz:

\begin{itemize}
  \item \code{ECONNREFUSED}: niemand hat die Verbindung angenommen. Typische Ursachen: falscher Port, Proxy nicht gestartet oder Upstream nicht erreichbar.
  \item \code{ECONNRESET}: die Gegenseite hat eine bestehende Verbindung zurueckgesetzt. Typische Ursache: der Upstream hat den Socket mitten im Stream geschlossen.
  \item \code{UND\_ERR\_SOCKET}: Nodes Undici-Client sah einen Socket-Fehler. Meist Reset oder kaputter Stream.
  \item \code{ETIMEDOUT}: Verbindung oder Antwort dauerte zu lange. Ursache: Netzwerkpfad oder Backend ist zu langsam.
  \item \code{AbortError}: wir haben den Request absichtlich abgebrochen. Typische Ursache: Claude Code hat die lokale Verbindung geschlossen.
\end{itemize}

Praktische Regel: Retry vor dem Senden der Response-Header; nach
Streaming-Beginn den Stream sauber mit einem Error-Event beenden. Sobald Header
und Teilinhalt gesendet wurden, kann der Proxy daraus keine normale JSON
\code{502}-Antwort mehr machen.

\section{Wichtige Dateien}

\subsection{\code{claude-o-go}}

Der Launcher exportiert API Key, Modell und lokale Base URL und startet dann
Claude Code:

\begin{verbatim}
export OPENCODE_GO_API_KEY="$api_key"
export OPENCODE_GO_MODEL="$model"
export ANTHROPIC_BASE_URL="http://127.0.0.1:${proxy_port}"
claude --bare --model "$claude_model_alias" "$@"
\end{verbatim}

Wichtige Defaults und Umgebungsvariablen:

\begin{itemize}
  \item \code{OPENCODE\_GO\_MODEL}: erforderlich; kein hartcodierter Launcher-Default.
  \item \code{OPENCODE\_GO\_PROXY\_PORT=4141}.
  \item \code{OPENCODE\_GO\_UPSTREAM\_RETRIES=2}.
  \item \code{CLAUDE\_CODE\_MODEL\_ALIAS=sonnet}.
  \item \code{CLAUDE\_CODE\_OPENCODE\_GO\_BARE=1}.
\end{itemize}

\subsection{\code{opencode-go-claude-proxy.mjs}}

Der zentrale Rewrite im Proxy:

\begin{verbatim}
body.model = model
\end{verbatim}

Bei Generierungsrequests streamt der Proxy die Antwort von OpenCode Go zurueck
an Claude Code. Er beachtet Node-Stream-Backpressure, bricht den
Upstream-Request ab, wenn Claude Code die lokale Verbindung schliesst,
wiederholt retrybare Upstream-Verbindungsfehler vor dem Senden der
Response-Header und wandelt Upstream-Resets mitten im Stream in ein sauberes
Streaming-Error-Event um, statt den Proxy-Prozess zu beenden.

Der Proxy implementiert ausserdem Endpunkte, die Claude Code abfragt:

\begin{verbatim}
HEAD /
HEAD /v1
GET /v1/models
POST /v1/messages/count_tokens
POST /v1/messages
\end{verbatim}

\section{Lokale Checks}

Vor dem Push ausfuehren:

\begin{verbatim}
npm run check
\end{verbatim}

Das prueft die Syntax beider JavaScript-Einstiegspunkte.

\section{Verifikation}

Starte:

\begin{verbatim}
OPENCODE_GO_PROXY_LOG=1 ./claude-o-go -p "Reply with exactly: OK"
\end{verbatim}

Proxy-Request-Logs sind im normalen Betrieb ausgeblendet, damit Claude Codes
TUI sauber bleibt. Mit \code{OPENCODE\_GO\_PROXY\_LOG=1} achte auf Zeilen wie:

\begin{verbatim}
OpenCode Go Claude proxy listening on http://127.0.0.1:4141/v1 -> qwen3.7-plus
POST /v1/messages
POST /messages -> 200 as qwen3.7-plus
OK
\end{verbatim}

Die wichtigste Zeile ist:

\begin{verbatim}
POST /messages -> 200 as qwen3.7-plus
\end{verbatim}

Das bedeutet: Claude Code hat an den lokalen Proxy gesendet, der Proxy hat an
OpenCode Go weitergeleitet, OpenCode Go hat HTTP 200 geliefert, und das Log
zeigt das verwendete Upstream-Modell.

\section{Haeufige Probleme}

\subsection{Proxy-Logs erscheinen in Claude Codes TUI}

Normale Request-Logs sind standardmaessig ausgeblendet. Zum Debuggen:

\begin{verbatim}
OPENCODE_GO_PROXY_LOG=1 claude-o-go
\end{verbatim}

Im normalen Betrieb unset lassen, damit die TUI sauber bleibt.

\subsection{Selected model issue}

Nutze den Launcher:

\begin{verbatim}
./claude-o-go
\end{verbatim}

Rufe \code{claude --model qwen3.7-plus} nicht direkt auf.

\subsection{Doppeltes \code{/v1/v1/messages}}

\code{ANTHROPIC\_BASE\_URL} muss auf die Proxy-Wurzel zeigen:

\begin{verbatim}
http://127.0.0.1:4141
\end{verbatim}

Sie darf nicht \code{/v1} enthalten, weil Claude Code selbst
\code{/v1/messages} anhaengt.

\subsection{Claude Code zeigt \code{sonnet}}

Das ist erwartet. Claude Code sieht \code{sonnet}; OpenCode Go erhaelt
\code{qwen3.7-plus}. Das Proxy-Log ist die Quelle der Wahrheit.

\section{Mental Model}

\begin{verbatim}
Claude Code CLI
  denkt model = sonnet
  sendet Anthropic Request
        |
        v
Lokaler Proxy
  schreibt model = qwen3.7-plus
  leitet weiter
        |
        v
OpenCode Go
  fuehrt das echte Modell aus
  gibt Antwort zurueck
        |
        v
Claude Code CLI
  zeigt die Antwort an
\end{verbatim}

Die Claude-Code-Nutzung bleibt gleich, aber das Modell-Backend ist OpenCode Go.

\end{document}
